\documentclass[11pt]{letter}
\usepackage[margin=1in]{geometry}
\usepackage{hyperref}
\signature{Jin Ho Choi}
\address{SmartEAR AI \\ Daegu, South Korea \\ jinhochoi@smartear.co.kr}
\begin{document}
\begin{letter}{Dr.\ Ahmet M.\ Elbir \\ Associate Editor \\ IEEE Transactions on Signal Processing}
\opening{Dear Dr.\ Elbir,}

We submit a new manuscript, ``A Closed-Loop Subspace--RFS Framework for DOA
Estimation and Multi-Target Tracking,'' which is a substantially revised version of our earlier submission
T-SP-35123-2026 (``\ldots via COP-RFS''). In accordance with IEEE SPS Policy
6.16 on the handling of rejected papers, we wish to make clear that this is a
\emph{substantially different} manuscript: it has a new central thesis, two new
theoretical results that were absent from the prior version, corrected
derivations, and a completely re-run experimental evaluation. A detailed
point-by-point response to all prior reviewer comments accompanies this
submission.

The prior reviews converged on one central concern: the work read as an
\emph{integration} of established components (the authors' COP front-end and a
GM-PHD tracker) rather than a fundamentally new contribution. We have addressed
this at the root by changing what the paper is \emph{about}:

\begin{itemize}
\item \textbf{The contribution is now the closed-loop coupling itself,
  established as a front-end-agnostic principle.} We prove (i) a posterior
  Cram\'er--Rao bound for the complete closed-loop estimator--tracker
  (Theorem~1), showing the temporal prior restores identifiability exactly in
  the underdetermined regime where the data-only bound diverges; and (ii) that
  the tracker-aided subspace refinement is minimum-variance on the Grassmann
  manifold under an explicit, gate-enforced condition (Theorem~2). These are new
  results, not present in the rejected version.
\item \textbf{We demonstrate the principle on two different front-ends} (the
  higher-order-cumulant COP and the covariance-based MUSIC), showing the closed
  loop benefits both---lowering localization error by $56\%$ (COP) and $44\%$
  (MUSIC) at low SNR---evidence that the contribution is general rather than a
  specific algorithm pairing.
\item \textbf{Every flagged derivation has been corrected}: the effective sample
  size (Kish), the basis-invariant subspace refinement (replacing the
  gauge-dependent interpolation), the velocity gate (removing double-counting),
  the cumulant-model assumptions, and the cumulant-noise covariance. The
  unsupported ``exact'' CRLB and the heuristic convergence bound have been
  removed; the back-end is renamed a track-oriented PHD with an honest statement
  of its relation to the exact recursion.
\item \textbf{The experiments are re-run with up to $500$ Monte-Carlo trials and
  $95\%$ confidence intervals}, with a closed-loop-MUSIC baseline, a
  $2\times2$ (front-end $\times$ loop) ablation, and a tracking back-end ablation
  against the LMB and $\delta$-GLMB labeled-RFS filters. We now explicitly \emph{delineate the operating regime}
  in which the higher-order front-end is advantageous, and we have removed all
  overclaims (``tracks all,'' ``substantially outperforms'').
\end{itemize}

Crucially, the revised paper now \emph{quantifies} the value of closing the
loop---a concrete gain in every regime where the temporal prior is informative
(restored identifiability, $-56\%$ low-SNR localization error, $-15\%$ identity
switches under target motion), with a provable no-harm guarantee elsewhere.
We believe these changes squarely address the reviewers' concerns about novelty
positioning, theoretical rigor, and experimental support, and that the resulting
paper makes a clear theoretical contribution appropriate for the Transactions.
Thank you for considering our work.

\closing{Sincerely,}
\end{letter}
\end{document}
